\scanpage{423}
\mspage{335}

En invoquant à nouveau 4.4.\unsure{2|1}, on trouve que $J\longrightarrow\unsure{M|\mathcal M}_{\pi/x}$ est cofinal, et comme ce foncteur se factorise par le foncteur d'inclusion \unsure{pl fidèle|pl. fidèle} $\underline S(\pi)_{/x}\hookrightarrow M_{\pi/x}$, il s'ensuit \deleted{que} comme tantôt que ce dernier est cofinal. Nous allons \deleted{\unsure{en conclure}} \inserted{admettre} pour un moment que ceci implique \unsure{que|pour} l'image inverse \deleted{$\unsure{J}$}$I_{\underline S}$ de cette sous-catégorie str. pleine, par le foncteur également cofinal $I\longrightarrow M_{\pi/x}$\unsure{;|,} est cofinal dans $I$. Ceci étant, l'équivalence $(**')$ (p. 333) implique $x\in\mathrm{Ob}\,\underline S$, ce qu'on voulait.

Donc il reste à prouver \deleted{\unsure{que le}} que si
\diagram{\begin{tikzcd}[ampersand replacement=\&]
 J'=\underline S(\pi)_{/x}\arrow[r,hook] \& M_{\pi/x}\unsure{=|\subset}J
\end{tikzcd}}
est cofinal, alors $I'=I_{\unsure{S|\underline S}}\longrightarrow I$ l'est aussi. Il semble qu'il faille \inserted{bel et bien} \unsure{(à nouveau|à nouveau} \deleted{\unsure{la}} invoquer les conditions de stabilité \deleted{\unsure{de}} $(\pi_0,\pi)$ \unsure{faite|faites} dans \unsure{d')|a')} (ce qui n'avait pas été nécessaire dans le cas c)), \unsure{et}

% archive sheet number: 422

